% BHU PYQ -- Physics Thermal Physics (End Term) (2024-25 NEP)
\documentclass[11pt,a4paper]{article}
\usepackage[a4paper,top=2.5cm,bottom=2.5cm,left=2.8cm,right=2.8cm,headheight=14pt]{geometry}
\usepackage{amsmath,amssymb}
\usepackage[T1]{fontenc}
\usepackage[utf8]{inputenc}
\usepackage{lmodern,microtype,fancyhdr,enumitem,hyperref,lastpage,parskip}

\hypersetup{colorlinks=true,urlcolor=black,linkcolor=black,pdftitle=PHYMJ21: Thermal Physics (End Term) -- BHU NEP 2024-25}

\pagestyle{fancy}\fancyhf{}
\renewcommand{\headrulewidth}{0.4pt}\renewcommand{\footrulewidth}{0.4pt}
\fancyhead[L]{\small   \textit{Banaras Hindu University}}
\fancyhead[R]{\small   \textit{PHYMJ21: Thermal Physics (End Term) (2024-25)}}
\fancyfoot[C]{\small Page  \thepage\ of \pageref{LastPage}}


\newlist{parts}{enumerate}{2}
\setlist[parts,1]{label=(\alph*),leftmargin=*,itemsep=4pt,topsep=3pt}
\setlist[parts,2]{label=(\roman*),leftmargin=*,itemsep=2pt,topsep=2pt}

\newcommand{\pts}[1]{\hfill   \textit{[#1]}}


\begin{document}

\begin{center}
  {\large\bfseries BANARAS HINDU UNIVERSITY}\\[2pt]
  {
\normalsize U.G. Semester 2 (NEP) Examination, 2024-25}\\[6pt]
  \rule{0.8\linewidth}{0.5pt}\\[6pt]
  {\large\bfseries PHYSICS}\\[4pt]
  {\large\bfseries Paper: PHYMJ21 --- Thermal Physics (End Term)}\\[4pt]
  {
\normalsize Time: 2:30 Hours \hfill Full Marks: 70}
\end{center}
\rule{\linewidth}{0.4pt}
\smallskip



\noindent   \textit{


\noindent   \textit{Instructions: Attempt   \textbf{FOUR} questions including Question No.~1, which is compulsory. The figures in the right-hand margin indicate marks.}
\bigskip




\noindent   \textbf{Question 1.}   \textit{(Compulsory --- Attempt any four)}

\begin{parts}
  \item State the Zeroth and First Laws of Thermodynamics.\pts{3.5}
  \item Derive the expression for efficiency of a Carnot engine: $\eta = 1 - \frac{T_2}{T_1}$.\pts{3.5}
  \item Define entropy. Show that the entropy of an isolated system never decreases.\pts{3.5}
  \item State the Clausius-Clapeyron latent heat equation and explain its physical significance.\pts{3.5}
\end{parts}

\bigskip



\noindent   \textbf{Question 2.}

\begin{parts}
  \item State and prove Maxwell's four thermodynamic relations.\pts{14}
\end{parts}

\bigskip



\noindent   \textbf{Question 3.}

\begin{parts}
  \item Derive Planck's radiation formula for blackbody radiation. Show how Wien's displacement law and Rayleigh-Jeans law can be deduced as limiting cases.\pts{14}
\end{parts}

\bigskip



\noindent   \textbf{Question 4.}

\begin{parts}
  \item Explain Joule-Thomson expansion. Derive an expression for the Joule-Thomson coefficient $\mu_{JT} = \left(\frac{\partial T}{\partial P}\right)_H$.\pts{14}
\end{parts}


\vfill\rule{\linewidth}{0.4pt}
\end{document}
